\documentclass[11pt]{article}

\usepackage[margin=1in]{geometry}
\usepackage{amsmath,amssymb,amsthm,mathtools}
\usepackage{graphicx}
\usepackage{microtype}
\usepackage[colorlinks=true,linkcolor=blue,citecolor=blue,urlcolor=blue]{hyperref}

\newtheorem{theorem}{Theorem}
\newtheorem{lemma}[theorem]{Lemma}
\newtheorem{remark}[theorem]{Remark}
\newcommand{\R}{\mathbb{R}}
\newcommand{\cH}{\mathcal{H}}
\newcommand{\eps}{\varepsilon}
\newcommand{\1}{\mathbf{1}}
\newcommand{\op}{\mathrm{op}}

\title{Linear Convergence for General Continuous Costs\\
in Quadratically Regularized Optimal Transport}
\author{Candidate solution packet for arXiv:2509.08547}
\date{9 August 2026}

\begin{document}
\maketitle

\begin{center}
\fbox{\parbox{0.9\textwidth}{\textbf{Status.}
Candidate full solution, likely valid pending expert review.  The result proves
the general-continuous-cost extension proposed in Remark~2.5 of the source
paper, under its original weak marginal hypotheses.}}
\end{center}

\section{Source question}

Gonz\'alez-Sanz, Nutz, and Riveros Valdevenito prove eventual linear
$L^2$ convergence of dual gradient descent for quadratically regularized
optimal transport with squared Euclidean cost \cite{GNR2026}.  In Remark~2.5
(page 7), they explain that the special cost is used only to show that certain
zero-level sections are negligible and ask whether their theorem extends to
general continuous costs.  The theorem and remark are reproduced below.

\begin{center}
\includegraphics[width=0.74\textwidth]{figures/open_problem_crop.png}
\end{center}

The result below answers that question affirmatively.  The new argument does
not prove convergence of the secant operators to a single linearization;
instead, it finds a common spectral gap for all eventual secant operators.

\section{Setup}

Let $P,Q$ be probability measures with compact topological supports
$\Omega,\Omega'$.  Assume that at least one support is connected; by symmetry
we take $\Omega'$ connected.  Let
$c:\Omega\times\Omega'\to\R$ be continuous and let $\eps>0$.  The QOT dual
objective is
\[
 \Gamma(f,g)=\int f\,dP+\int g\,dQ
 -\frac1{2\eps}\int (f(x)+g(y)-c(x,y))_+^2\,dP(x)dQ(y).
\]
We use the normalized Hilbert space
\[
 \cH_0=\left\{(f,g)\in L^2(P)\oplus L^2(Q):
                  \int f\,dP=\int g\,dQ\right\}.
\]
For continuous initial data $(f_0,g_0)\in\cH_0$ and step size
$0<\eta<\eps$, define
\[
 (f_{n+1},g_{n+1})=(f_n,g_n)+\eta\,D\Gamma(f_n,g_n).
\]
The gradient belongs to $\cH_0$, so normalization is preserved.  Let
$(f_*,g_*)$ be the normalized optimal potentials and set
\[
 \xi_n(x,y)=f_n(x)+g_n(y)-c(x,y),\qquad
 \xi_*(x,y)=f_*(x)+g_*(y)-c(x,y).
\]
The first-order equations are
\begin{equation}
\label{eq:FOC}
 \int (\xi_*(x,y))_+\,dQ(y)=\eps\quad(x\in\Omega),
 \qquad
 \int (\xi_*(x,y))_+\,dP(x)=\eps\quad(y\in\Omega').
\end{equation}

\section{Proof intuition}

The derivative of the positive part between $\xi_n$ and $\xi_*$ is an exact
secant weight $w_n(x,y)\in[0,1]$.  A large zero-level set can make these
weights oscillate and prevents the operator-norm limit used in the source.
But no limit is needed.

Equation~\eqref{eq:FOC} forces a fixed strictly active core
$E=\{\xi_*>\eps/2\}$ to have uniformly positive sections.  On the normalized
space, the corresponding energy
\[
 \int_E(f(x)+g(y))^2\,dP(x)dQ(y)
\]
has a positive spectral gap.  Connectedness is exactly what eliminates its
kernel: a zero-energy pair is locally a pair of opposite constants along the
open bipartite core, hence globally so, and normalization removes the
constant.

After burn-in, both $\xi_n$ and $\xi_*$ are positive on $E$, so $w_n=1$ there.
Every exact secant energy therefore dominates the same gapped core.  Its upper
spectral bound is universally $2$, and $0<\eta/\eps<1$ turns the exact error
recursion into a uniform strict contraction.

\section{Main theorem}

\begin{theorem}[General continuous costs]
\label{thm:main}
Let $P,Q$ have compact supports and suppose that at least one support is
connected.  Let $c$ be continuous, $(f_0,g_0)\in\cH_0$ be continuous, and
$0<\eta<\eps$.  Then there exist $n_0\in\mathbb N$, $C<\infty$, and
$\delta\in(0,1)$ such that
\[
 \|(f_n,g_n)-(f_*,g_*)\|_{L^2(P)\oplus L^2(Q)}
 \le C\delta^n\qquad(n\ge n_0).
\]
Thus Theorem~2.3 of \cite{GNR2026} extends to every continuous transport cost.
In particular it holds under all the marginal assumptions of that theorem;
the condition that $P,Q$ avoid boundaries of convex sets is unnecessary for
this extension.
\end{theorem}

\section{Proof}

\subsection{Uniform convergence before linearization}

\begin{lemma}
\label{lem:preconvergence}
In the setting of Theorem~\ref{thm:main},
\[
 \|\xi_n-\xi_*\|_{L^\infty(\Omega\times\Omega')}\longrightarrow0.
\]
\end{lemma}

\begin{proof}
The source paper proves this before its only use of the squared Euclidean
form; we record why the argument needs only continuity of $c$.

First, the elementary positive-part inequalities and $\eta/\eps\le1$ show
that the sup-norm distance of the iterates from the potentials is uniformly
bounded (up to the harmless potential translation).  This is the proof of
Lemma~5.2 in \cite{GNR2026} and is independent of $c$.

Second, the iterates have a common modulus of continuity.  Indeed, for two
points $x,x'$, write the secant identity
\[
 (u)_+-(v)_+=\theta(u-v),\qquad 0\le\theta\le1,
\]
inside the $Q$-integral defining the update of $f_n$.  If $\rho_n$ is a
modulus for $f_n$ and
\[
 \omega_c(r)=\sup_{\substack{|x-x'|\le r\\y\in\Omega'}}
              |c(x,y)-c(x',y)|,
\]
then the difference of the two updates is a convex combination, with total
weight at most $\eta/\eps$, of a term bounded by $\rho_n(r)$ and one bounded
by $\omega_c(r)$.  Hence $f_{n+1}$ has modulus at most
$\max(\rho_n,\omega_c)$.  The same holds for $g_n$.  Starting from continuous
data gives a common modulus for all iterates.

The map $D\Gamma$ is $2/\eps$-Lipschitz on the Hilbert space, again by the
one-Lipschitz property of the positive part.  Since $0<\eta<\eps$, the usual
ascent inequality gives
\[
 \sum_{n=0}^\infty\|D\Gamma(f_n,g_n)\|_2^2<\infty.
\]
Uniform boundedness and equicontinuity give compactness by Arzel\`a--Ascoli.
Every cluster point is stationary.  Existence and uniqueness of continuous
potentials modulo translation for a continuous QOT cost, with uniqueness
under connectedness of one support, are the general-cost preliminary results
used in \cite{GNR2026}; their proof uses \eqref{eq:FOC} and continuity only.
Normalization therefore makes $(f_*,g_*)$ the unique cluster point.  Thus
$f_n\oplus g_n\to f_*\oplus g_*$ uniformly, which is precisely the asserted
convergence of $\xi_n$.
\end{proof}

\subsection{A fixed active core has a spectral gap}

Let $B=\|(\xi_*)_+\|_\infty$.  Equation~\eqref{eq:FOC} implies $B\ge\eps$.
Set
\[
 \tau=\frac\eps2,\qquad
 E=\{(x,y):\xi_*(x,y)>\tau\},\qquad
 \lambda=\frac{\eps}{2B}>0.
\]
For all $x\in\Omega$ and $y\in\Omega'$, respectively,
\begin{equation}
\label{eq:sectionmass}
 Q(E_x)\ge\lambda,\qquad P(E^y)\ge\lambda.
\end{equation}
Indeed,
\[
 \eps=\int(\xi_*)_+\,dQ
 \le \tau+BQ\{\xi_*>\tau\},
\]
and the other inequality is symmetric.

For $(f,g)\in\cH_0$, define the core form
\[
 \mathfrak a_E(f,g)=\int_E(f(x)+g(y))^2\,dP(x)dQ(y).
\]

\begin{lemma}[Active-core gap]
\label{lem:gap}
There is a constant $\kappa>0$ such that
\[
 \mathfrak a_E(f,g)\ge
 \kappa\bigl(\|f\|_{L^2(P)}^2+\|g\|_{L^2(Q)}^2\bigr)
 \qquad((f,g)\in\cH_0).
\]
\end{lemma}

\begin{proof}
Let $q(x)=Q(E_x)$ and $p(y)=P(E^y)$.  The self-adjoint operator represented by
$\mathfrak a_E$ has diagonal multiplication part
$D=\operatorname{diag}(q,p)$ and cross operator
\[
 Tg(x)=\int\1_E(x,y)g(y)\,dQ(y),\qquad T^*f(y)=\int\1_E(x,y)f(x)\,dP(x).
\]
The kernel $\1_E$ belongs to $L^2(P\otimes Q)$, so $T$ and $T^*$ are
Hilbert--Schmidt.  By \eqref{eq:sectionmass}, the compression of $D$ to
$\cH_0$ is bounded below by $\lambda$.  The compressed core operator is
therefore an invertible positive operator plus a compact operator, hence is
Fredholm.

It remains to show that its kernel is zero.  If $\mathfrak a_E(f,g)=0$, then
$f(x)+g(y)=0$ for $P\otimes Q$-almost every $(x,y)\in E$.  Fix
$y_0\in\Omega'$.  By \eqref{eq:sectionmass}, choose $x_0$ with
$(x_0,y_0)\in E$.  Since $E$ is relatively open, there are relative-open
neighborhoods $U\ni x_0$ and $V\ni y_0$ such that $U\times V\subset E$.
Both have positive measure because $\Omega,\Omega'$ are topological supports.
Fubini's theorem now shows that $g$ is $Q$-almost everywhere constant on $V$.

These neighborhoods cover the connected set $\Omega'$.  Constants agree on
every nonempty overlap because such an overlap is relative-open and has
positive $Q$-measure.  Connectedness then yields $g=c$ $Q$-almost everywhere
on $\Omega'$.  Since every $E_x$ has positive $Q$-measure, the zero-energy
identity gives $f=-c$ $P$-almost everywhere.  The normalization
$\int f\,dP=\int g\,dQ$ forces $c=0$.

Thus the compressed Fredholm operator has trivial kernel.  Being self-adjoint
of index zero, it is invertible; its positive quadratic form consequently has
some lower bound $\kappa>0$.
\end{proof}

\subsection{Every eventual secant operator dominates the core}

Let $e_n=(f_n-f_*,g_n-g_*)\in\cH_0$.  Define the measurable secant slope
\[
 w_n(x,y)=
 \begin{cases}
 \dfrac{(\xi_n(x,y))_+-(\xi_*(x,y))_+}
       {\xi_n(x,y)-\xi_*(x,y)},&\xi_n(x,y)\ne\xi_*(x,y),\\[1.2ex]
 \1_{\{\xi_*(x,y)>0\}},&\xi_n(x,y)=\xi_*(x,y).
 \end{cases}
\]
Then $0\le w_n\le1$, and the gradient update gives the exact recurrence
\begin{equation}
\label{eq:recurrence}
 e_{n+1}=\left(I-\frac\eta\eps A_n\right)e_n,
\end{equation}
where $A_n$ is the self-adjoint operator on $\cH_0$ with form
\[
 \langle A_n(f,g),(f,g)\rangle
 =\int w_n(x,y)(f(x)+g(y))^2\,dP(x)dQ(y).
\]

By Lemma~\ref{lem:preconvergence}, for all sufficiently large $n$ we have
$\|\xi_n-\xi_*\|_\infty<\tau$.  On $E$, both $\xi_*>\tau$ and $\xi_n>0$,
so $w_n=1$.  Lemma~\ref{lem:gap} then gives
\begin{equation}
\label{eq:lower}
 \langle A_nh,h\rangle\ge\kappa\|h\|^2\qquad(h\in\cH_0).
\end{equation}
For the upper bound, write
$m=\int f\,dP=\int g\,dQ$.  Since $w_n\le1$,
\begin{align*}
 \langle A_n(f,g),(f,g)\rangle
 &\le\int(f(x)+g(y))^2\,dP(x)dQ(y)\\
 &=\|f\|_2^2+\|g\|_2^2+2m^2\\
 &\le2(\|f\|_2^2+\|g\|_2^2).
\end{align*}
Thus the spectrum of every eventual $A_n$ is contained in $[\kappa,2]$.
Putting $a=\eta/\eps\in(0,1)$, the spectral theorem and
\eqref{eq:recurrence} yield
\[
 \|e_{n+1}\|\le\delta\|e_n\|,
 \qquad
 \delta=\max\{|1-a\kappa|,|1-2a|\}<1.
\]
Iteration proves the stated geometric estimate.
\qed

\section{Relation to later literature}

The later PL-inequality paper \cite{GNRPL2026} independently proves linear
convergence for arbitrary uniformly continuous costs and provides explicit
rates, but assumes that one marginal has a density bounded above and away
from zero on a convex support (or on the closure of a bounded Lipschitz
domain).  Theorem~\ref{thm:main} is non-quantitative but retains the source
paper's substantially weaker marginal framework; its active-core Fredholm
argument requires only compact supports and connectedness of one support.

\section{Verification, limitations, and novelty}

The proof was audited step by step in \path{VERIFICATION.md}.  The optional
matrix checker \path{code/verify_weighted_core.py} verifies the spectral
domination and contraction inequalities for finite weighted bipartite cores;
it is not part of the proof.

The theorem gives a qualitative contraction factor because the Fredholm gap
$\kappa$ is not explicit.  It concerns fixed-step dual gradient descent with
$0<\eta<\eps$ and does not assert convergence at or above the threshold.

A bounded literature search on 9 August 2026 used the exact source title and
arXiv id together with “general continuous cost,” “strictly active,” “secant
weights,” “spectral gap,” and “linear convergence.”  It found the source and
the stronger-assumption PL result \cite{GNRPL2026}, but no theorem with the
weak marginal hypotheses above and no use of the present secant-domination
argument.  Novelty confidence is moderate pending specialist review.

\section{Human-review recommendation}

The packet is ready for expert review as a candidate full solution of
Remark~2.5.  The principal checks are the local-a.e.-constancy patching in
Lemma~\ref{lem:gap}, the Fredholm compression to $\cH_0$, and the claim that
the source's preconvergence proof remains valid for every continuous cost.

\begin{thebibliography}{9}

\bibitem{GNR2026}
A. Gonz\'alez-Sanz, M. Nutz, and A. Riveros Valdevenito,
\emph{Linear convergence of gradient descent for quadratically regularized
optimal transport}, arXiv:2509.08547v3, 2026.

\bibitem{GNRPL2026}
A. Gonz\'alez-Sanz, M. Nutz, and A. Riveros Valdevenito,
\emph{Polyak--\L{}ojasiewicz inequality for quadratically regularized optimal
transport}, arXiv:2605.27175, 2026.

\end{thebibliography}

\end{document}
